\section{Разработка алгоритма укладки паллеты}

В данном разделе описываются алгоритмы укладки коробок на паллету, разработанные в рамках проекта.

\subsection{Структура HeightHandler}

Ключевым компонентом решения является структура \texttt{HeightHandler}, которая эффективно хранит информацию о текущих высотах на паллете. Данная структура критически важна для производительности алгоритма, поскольку операции с ней выполняются при размещении каждой коробки.

\subsubsection{Основные операции}

HeightHandler предоставляет следующие ключевые операции:

\begin{itemize}
    \item \texttt{get\_h(x, y, X, Y)} --- возвращает максимальную высоту в прямоугольной области $[x, X] \times [y, Y]$
    \item \texttt{get\_area(x, y, X, Y)} --- возвращает площадь (в мм²), которая находится на максимальной высоте в заданном прямоугольнике (используется для расчёта опоры коробки)
    \item \texttt{add\_rect(x, y, X, Y, h)} --- добавляет новый прямоугольник с заданной высотой, обновляя карту высот
    \item \texttt{get\_dots(header, box)} --- возвращает ``интересные'' точки для размещения коробки заданного размера
\end{itemize}

Функция \texttt{get\_dots} генерирует кандидатные позиции на основе углов существующих прямоугольников, что значительно сокращает пространство поиска с $O(W \times L)$ до $O(n)$, где $n$ --- количество уже уложенных коробок.

\subsubsection{Реализации HeightHandler}

В рамках работы были исследованы пять различных реализаций структуры HeightHandler с разными алгоритмическими подходами.

\paragraph{Baseline (наивная реализация)}

Простейшая реализация, которая хранит все добавленные прямоугольники в векторе и при каждом запросе честно перебирает их.

\textbf{Сложность операций:}
\begin{itemize}
    \item \texttt{get\_h()}: $O(n)$ --- перебор всех прямоугольников
    \item \texttt{get\_area()}: $O(n)$ --- перебор всех прямоугольников
    \item \texttt{add\_rect()}: $O(1)$ --- просто добавление в конец вектора
\end{itemize}

\textbf{Особенность:} Корректность гарантируется свойством задачи паллетизации --- прямоугольники добавляются ``сверху'', поэтому новый прямоугольник всегда имеет высоту строго больше, чем все существующие в той же области. Два прямоугольника на одной высоте никогда не перекрываются.

\paragraph{Rects (оптимизированный список прямоугольников)}

Улучшенная версия Baseline с оптимизацией: прямоугольники сортируются по убыванию высоты. Это позволяет при поиске максимума остановиться на первом найденном пересечении.

\textbf{Сложность операций:}
\begin{itemize}
    \item \texttt{get\_h()}: $O(n)$ в худшем случае, но на практике значительно быстрее из-за раннего выхода
    \item \texttt{get\_area()}: $O(n)$ в худшем случае, раннее завершение при снижении высоты
    \item \texttt{add\_rect()}: $O(n)$ --- вставка с сохранением сортировки
\end{itemize}

\paragraph{Grid (плотная 2D сетка)}

Разбивает паллету на ячейки размером $10 \times 10$ мм. В каждой ячейке хранится максимальная высота. Для паллеты $1200 \times 800$ мм это $120 \times 80 = 9600$ ячеек.

\textbf{Сложность операций:}
\begin{itemize}
    \item \texttt{get\_h()}: $O(k)$, где $k$ --- количество ячеек в запросе
    \item \texttt{get\_area()}: $O(k)$
    \item \texttt{add\_rect()}: $O(k)$ --- обновление всех затронутых ячеек
\end{itemize}

\textbf{Недостаток:} Большой объём памяти и медленное обновление при добавлении крупных прямоугольников.

\paragraph{SegTree (2D дерево отрезков)}

Двумерное дерево отрезков с отложенным обновлением (Lazy Propagation). X-дерево содержит Y-деревья, каждый Y-узел хранит количество 2D-ячеек с максимальной высотой.

\textbf{Сложность операций:}
\begin{itemize}
    \item \texttt{get\_h()}: $O(\log W \cdot \log H)$
    \item \texttt{get\_area()}: $O(\log W \cdot \log H)$
    \item \texttt{add\_rect()}: $O(\log W \cdot H)$ --- необходим merge Y-деревьев
\end{itemize}

\textbf{Особенность:} Асимптотически эффективен для запросов, но высокая константа и сложная операция обновления делают его медленнее на практике.

\paragraph{Quadtree (квадродерево)}

Рекурсивное разбиение пространства на 4 части. Каждый узел хранит максимальную высоту в своей области. При частичном пересечении с областью обновления узел разбивается на 4 дочерних.

\textbf{Сложность операций:}
\begin{itemize}
    \item \texttt{get\_h()}: $O(\log(\max(W,H)))$ при равномерном распределении
    \item \texttt{get\_area()}: $O(n)$ --- использует fallback на список прямоугольников
    \item \texttt{add\_rect()}: $O(\log(\max(W,H)))$ в среднем, но может деградировать
\end{itemize}

\textbf{Недостаток:} Высокие накладные расходы на управление памятью (unique\_ptr), частое разбиение узлов.

\subsubsection{Сравнительное тестирование}

Тестирование проводилось на 436 реальных задачах с суммарно 27.5 миллионами операций.

\begin{table}[ht]
\centering
\caption{Результаты тестирования реализаций HeightHandler (время в мс)}
\label{tab:heighthandler}
\begin{tabular}{|l|r|r|r|r|r|}
\hline
\textbf{Алгоритм} & \textbf{Total} & \textbf{add\_rect} & \textbf{get\_h} & \textbf{get\_area} & \textbf{Speedup} \\
\hline
Baseline & 36412.1 & 3.9 & 19724.4 & 21230.7 & 1.00x \\
Rects & 1657.5 & 13.7 & 1713.4 & 2121.1 & \textbf{21.97x} \\
Grid & 120688.0 & 111.6 & 32015.1 & 93546.9 & 0.30x \\
SegTree & 71199.2 & 5824.6 & 34876.8 & 33292.1 & 0.51x \\
Quadtree & 168970.6 & 885.0 & 86821.8 & 81856.8 & 0.22x \\
\hline
\end{tabular}
\end{table}

\subsubsection{Анализ результатов тестирования HeightHandler}

Реализация \textbf{Rects} показала наилучшую производительность с ускорением в \textbf{21.97x} относительно базовой реализации. Это объясняется:
\begin{itemize}
    \item Сортировкой по убыванию высоты, позволяющей быстро находить максимум
    \item Низкими накладными расходами (простой вектор)
    \item Эффективным использованием кэша процессора
\end{itemize}

\textbf{Grid, SegTree и Quadtree} оказались медленнее Baseline. Причины:
\begin{itemize}
    \item Высокие накладные расходы на управление структурами данных
    \item Типичные коробки покрывают значительную часть паллеты, что нивелирует преимущества иерархических структур
    \item Константные факторы асимптотически эффективных алгоритмов превышают выигрыш
\end{itemize}

В финальной версии алгоритма используется реализация \textbf{Rects}.

\subsection{Жадный алгоритм (GreedySolver)}

GreedySolver реализует простой, но эффективный жадный подход к укладке коробок.

\begin{algorithm}[H]
\caption{GreedySolver}
\label{alg:greedy}
\KwIn{boxes --- list of boxes, pallet --- pallet dimensions}
\KwOut{answer --- positions of all boxes}
$answer \leftarrow$ empty list\;
$heightHandler \leftarrow$ new HeightHandler($pallet.length$, $pallet.width$)\;
\While{boxes is not empty}{
    $bestScore \leftarrow \infty$\;
    $bestBox \leftarrow$ null\;
    $bestPosition \leftarrow$ null\;
    $box \leftarrow boxes[0]$\;
    $availableRotations \leftarrow$ getAvailableBoxes($box$)\;
    \ForEach{rotatedBox $\in$ availableRotations}{
        $candidatePositions \leftarrow$ heightHandler.getDots($rotatedBox$)\;
        \ForEach{$(x, y) \in candidatePositions$}{
            $h \leftarrow$ heightHandler.getH($x$, $y$, $x + rotatedBox.length$, $y + rotatedBox.width$)\;
            $score \leftarrow h + rotatedBox.height$\;
            \If{$score < bestScore$}{
                $bestScore \leftarrow score$\;
                $bestBox \leftarrow rotatedBox$\;
                $bestPosition \leftarrow (x, y, h)$\;
            }
        }
    }
    answer.add($bestBox$ at $bestPosition$)\;
    heightHandler.addRect($bestPosition$, $bestBox.height$)\;
    remove $box$ from $boxes$\;
}
\Return{answer}
\end{algorithm}

\textbf{Сложность:} $O(n \cdot k \cdot m)$, где $n$ --- число коробок, $k$ --- число ориентаций (до 6), $m$ --- число кандидатных позиций.

\subsection{Жадный поиск в большой окрестности (LNSSolver)}

LNSSolver использует метаэвристический подход Large Neighborhood Search (LNS) с жадным критерием принятия решений.

\subsubsection{Представление решения}

Решение кодируется структурой \texttt{SimulationParams}:
\begin{itemize}
    \item \texttt{order} --- массив структур \texttt{BoxMeta}, определяющий порядок укладки
    \item \texttt{support\_threshold} --- минимальный порог опоры (коробки с меньшей опорой отклоняются)
\end{itemize}

Структура \texttt{BoxMeta}:
\begin{itemize}
    \item \texttt{box\_id} --- индекс типа коробки в исходном наборе
    \item \texttt{rotation} --- выбранная ориентация ($-1$ = рассмотреть все, $0$--$5$ = фиксированная)
    \item \texttt{position} --- предпочтительный угол паллеты ($0$--$3$, по умолчанию $3$)
\end{itemize}

\textbf{Порядок элементов в массиве определяет порядок укладки коробок.} Перестановка элементов приводит к другой укладке.

\subsubsection{Функция симуляции}

Функция \texttt{simulate} принимает \texttt{SimulationParams} и выполняет укладку:
\begin{enumerate}
    \item Для каждой коробки в порядке \texttt{order}:
    \begin{enumerate}
        \item Если \texttt{rotation} $\neq -1$ --- используем указанную ориентацию, иначе перебираем все
        \item Для каждой ориентации получаем кандидатные позиции из HeightHandler
        \item Фильтруем позиции: отклоняем, если центр масс коробки не попадает на опору
        \item Фильтруем позиции: отклоняем, если \texttt{support\_ratio} $<$ \texttt{support\_threshold}
        \item Выбираем позицию с минимальной итоговой высотой $h + h_{box}$
        \item При равной высоте выбираем позицию ближе к углу \texttt{position}
    \end{enumerate}
    \item Возвращается укладка и метрики (percolation, min\_support\_ratio, center\_of\_mass)
\end{enumerate}

\subsubsection{Операторы мутации}

Реализовано 13 операторов мутации, разделённых на категории:

\textbf{Изменение параметров (position/rotation):}
\begin{itemize}
    \item \texttt{position\_rotation\_fixed} --- установка одинакового значения для сегмента
    \item \texttt{position\_rotation\_random} --- случайные значения для каждой коробки в сегменте
\end{itemize}

\textbf{Изменение порядка:}
\begin{itemize}
    \item \texttt{reverse} --- реверс подмассива
    \item \texttt{shuffle} --- случайная перестановка подмассива
    \item \texttt{swap} --- обмен двух случайных коробок
    \item \texttt{swap\_adjacent} --- обмен соседних коробок
\end{itemize}

\textbf{Сортировки:}
\begin{itemize}
    \item \texttt{sort\_by\_area} --- сортировка по площади основания
    \item \texttt{sort\_by\_volume} --- сортировка по объёму
    \item \texttt{sort\_by\_height} --- сортировка по высоте
    \item \texttt{sort\_by\_weight} --- сортировка по весу
\end{itemize}

\textbf{Специальные:}
\begin{itemize}
    \item \texttt{threshold} --- изменение порога минимальной опоры
    \item \texttt{group\_by\_sku} --- группировка коробок одного типа вместе
    \item \texttt{move\_bad\_box} --- перемещение коробки с худшей опорой в начало
\end{itemize}

Каждому оператору назначается вес, определяющий вероятность его выбора.

\begin{algorithm}[H]
\caption{LNSSolver}
\label{alg:lns}
\KwIn{testData --- input data, endTime --- time limit}
\KwOut{bestAnswer --- best found solution}
$params \leftarrow$ initializeParams($testData$)\;
$bestResult \leftarrow$ simulate($testData$, $params$)\;
$bestParams \leftarrow params$\;
\While{$currentTime < endTime$}{
    $newParams \leftarrow$ copy($params$)\;
    $mutation \leftarrow$ selectMutation($weights$)\;
    newParams.mutate($mutation$)\;
    $newResult \leftarrow$ simulate($testData$, $newParams$)\;
    \If{$newResult.score > bestResult.score$}{
        $bestResult \leftarrow newResult$\;
        $bestParams \leftarrow newParams$\;
        $params \leftarrow newParams$\;
    }
}
\Return{bestResult.answer}
\end{algorithm}

\subsubsection{Целевая функция и критерий принятия}

Целевая функция (score) вычисляется как взвешенная сумма метрик:

$$score = w_1 \cdot percolation + w_2 \cdot min\_support\_ratio + w_3 \cdot (1 - com\_z\_normalized)$$

где:
\begin{itemize}
    \item $w_1$, $w_2$, $w_3$ --- настраиваемые веса (по умолчанию $w_1 = 1$, $w_2 = 1$, $w_3 = 10$)
    \item $com\_z\_normalized$ --- нормализованная высота центра масс
\end{itemize}

Алгоритм использует \textbf{жадный критерий принятия}: новое решение принимается только если $score_{new} > score_{best}$.

\subsection{Метрики качества укладки}

Для оценки качества укладки используются следующие метрики:

\begin{itemize}
    \item \textbf{Percolation} (коэффициент заполнения) --- отношение суммарного объёма коробок к объёму паллеты:
    $$percolation = \frac{V_{boxes}}{L \times W \times H}$$
    
    \item \textbf{Min support ratio} --- минимальный коэффициент опоры среди всех коробок. Для каждой коробки рассчитывается как отношение площади, опирающейся на предыдущий слой, к полной площади основания.
    
    \item \textbf{Center of mass} --- координаты центра масс паллеты, используются для оценки устойчивости.
\end{itemize}

\subsection{Результаты}

\begin{table}[ht]
\centering
\caption{Сравнительные результаты алгоритмов}
\label{tab:benchmark_eyu}
\begin{tabular}{|l|c|c|c|c|c|}
\hline
\textbf{Алгоритм} & \textbf{Лимит} & \textbf{Среднее} & \textbf{Мин} & \textbf{Макс} & \textbf{Время} \\
\hline
GreedySolver & — & 73.56\% & 59.01\% & 85.21\% & 308 мс \\
LNSSolver & 1 сек & 77.18\% & 72.86\% & 86.39\% & 14 сек \\
LNSSolver & 5 сек & 78.23\% & 73.11\% & 88.15\% & 70 сек \\
LNSSolver & 30 сек & 78.93\% & 75.36\% & 88.65\% & 7 мин \\
LNSSolver & 300 сек & 79.71\% & 76.03\% & 90.39\% & 70 мин \\
\hline
\end{tabular}
\end{table}

\subsubsection{Анализ результатов}

\textbf{GreedySolver} демонстрирует высокую скорость работы --- все 436 тестов обрабатываются за 308 миллисекунд. Однако качество укладки варьируется в широком диапазоне (59\%--85\%).

\textbf{LNSSolver} показывает значительно более стабильные результаты: минимальный коэффициент заполнения составляет 72.86\% даже при лимите в 1 секунду, что выше худшего результата GreedySolver.

Наблюдается логарифмическая зависимость качества от времени работы:
\begin{itemize}
    \item Первая секунда даёт прирост +3.6 п.п. относительно GreedySolver
    \item Увеличение времени в 5 раз (1с $\rightarrow$ 5с) даёт +1.05 п.п.
    \item Увеличение времени в 6 раз (5с $\rightarrow$ 30с) даёт +0.7 п.п.
    \item Увеличение времени в 10 раз (30с $\rightarrow$ 300с) даёт +0.78 п.п.
\end{itemize}

\subsubsection{Практические рекомендации}

\begin{enumerate}
    \item Для систем реального времени рекомендуется использовать LNSSolver с лимитом 1--5 секунд
    \item Для офлайн-планирования оптимален лимит 30--300 секунд
    \item GreedySolver может использоваться для быстрой оценки или как начальное приближение
\end{enumerate}

\subsection{Выводы}

\begin{enumerate}
    \item Разработанный алгоритм LNSSolver достигает коэффициента заполнения до 90\%, что соответствует современным промышленным решениям
    \item Жадный поиск в большой окрестности (LNS) с разнообразными операторами мутации эффективно исследует пространство решений, обеспечивая стабильное качество на всех тестах
    \item Структура HeightHandler обеспечивает эффективную работу с картой высот, что критически важно для производительности
    \item Простая реализация HeightHandler на списке прямоугольников показала лучшие результаты по сравнению с более сложными структурами данных (деревья отрезков, квадродеревья)
\end{enumerate}
